\section{CONCLUSION}
\label{sec:conclusion}

This paper framed layout-preserving PDF translation as a three-stage document
problem rather than a translation problem with a rendering tail. The implemented
pipeline analyses pages visually, translates addressable elements under an
identifier contract, and reconstructs the PDF with font-size and colour
decisions computed from the source page.

The stage-level evidence supports that framing. On OmniDocBench, the parser
recovers regions well under a granularity-invariant matcher, labels them
reliably, and preserves reading order; its harmful failures are the regions,
formulas, tables and OCR slices that never reach translation correctly. On
WMT24++, competitive LLMs differ less by model choice than target languages
differ from one another. Finally, a prompt-level $\pm15\%$ character-length
constraint is weakly obeyed and uncorrelated with adequacy, so layout fidelity
must be enforced by reconstruction rather than by the translator.

The work leaves one main task: evaluate complete translated PDFs on a shared
real-document benchmark, with metrics for adequacy, overflow, readable font
size, colour fidelity, visual similarity, table-cell accuracy and runtime. That
benchmark would show whether the stage-level bottleneck identified here becomes
the dominant end-to-end bottleneck in practice.
